% ===== Document Preamble =====
\documentclass[11pt]{article}
\usepackage[margin=0.5in]{geometry}
\usepackage{amsmath,amssymb}
\usepackage{xcolor}
\usepackage{enumitem}
\usepackage{needspace}
\usepackage{comment}

% ===== Colors =====
\definecolor{AnsRed}{RGB}{170,0,0}
\definecolor{WrongBlue}{RGB}{0,0,160}
\definecolor{SkipGray}{RGB}{120,120,120}

% ===== Choice list styling (tight, consistent) =====
\setlist[enumerate]{%
  leftmargin=1.5em,
  itemsep=0.25em,
  topsep=0.25em,
  label=\textbf{\Alph*.}
}

% ===== Style Wrappers =====
\newcommand{\QuestionStyle}[1]{{\large\bfseries\boldmath #1}}
\newcommand{\AnswerStyle}[1]{{\large\bfseries\boldmath #1}}
\newcommand{\ReasonStyle}[1]{{\normalsize #1}}

% ===== Show / Hide Reasoning Toggle =====
\newif\ifshowreason
\showreasontrue      % Instructor version
% \showreasonfalse   % Student version

% ===== Choice Item Macros (Answer + Reasoning) =====
% Usage:
%   \Correct{Answer}{Reasoning}
%   \Wrong{Answer}{Reasoning}

\newcommand{\Correct}[2]{%
  \item \AnswerStyle{\textcolor{AnsRed}{#1}}%
  \ifshowreason\\[-0.2em]
    \ReasonStyle{\textcolor{AnsRed}{#2}}%
  \fi
}

\newcommand{\Wrong}[2]{%
  \item \AnswerStyle{\textcolor{WrongBlue}{#1}}%
  \ifshowreason\\[-0.2em]
    \ReasonStyle{\textcolor{WrongBlue}{#2}}%
  \fi
}

% ===== Multiple-Choice Question Block Macro =====
% Usage: \QBlock{Q1}{stem...}{<items>}{answer key}

\newcommand{\QBlock}[4]{%
  \Needspace{8\baselineskip}%
  \noindent
  \begin{minipage}{\linewidth}
    \QuestionStyle{#1.~#2}

    \vspace{0.25em}
    \begin{enumerate}
      #3
    \end{enumerate}

    \vspace{0.1em}
    \noindent\textit{\textbf{Answer:}~#4}

    \vspace{0.25em}
    \hrule height 0.4pt
  \end{minipage}
  \vspace{0.75em}
}

% ===== Open-Ended Question Block Macro =====
% Usage: \OEBlock{Q1}{stem...}{answer text}

\newcommand{\OEBlock}[3]{%
  \Needspace{5\baselineskip}%
  \noindent
  \begin{minipage}{\linewidth}
    \QuestionStyle{#1.~#2}

    \vspace{0.25em}
    \noindent\textit{\textbf{Answer:}~#3}

    \vspace{0.25em}
    \hrule height 0.4pt
  \end{minipage}
  \vspace{0.75em}
}

% ===== Multi-Part Question Block Macro =====
% Usage: \MPBlock{Q1}{stem...}{<parts>}

\newcommand{\MPBlock}[3]{%
  \Needspace{8\baselineskip}%
  \noindent
  \begin{minipage}{\linewidth}
    \QuestionStyle{#1.~#2}

    \vspace{0.25em}
    #3

    \vspace{0.25em}
    \hrule height 0.4pt
  \end{minipage}
  \vspace{0.75em}
}

% ===== MC Sub-Part Macro =====
% Usage: \MCPart{a}{stem...}{<items>}{answer key}

\newcommand{\MCPart}[4]{%
  \par\vspace{0.35em}
  \noindent\textbf{(#1)}~#2
  \begin{enumerate}
    #3
  \end{enumerate}
  \vspace{0.1em}
  \noindent\textit{\textbf{Answer:}~#4}
  \par
}

% ===== Open-Ended Sub-Part Macro =====
% Usage: \OEPart{a}{stem...}{answer text}

\newcommand{\OEPart}[3]{%
  \par\vspace{0.35em}
  \noindent\textbf{(#1)}~#2
  \par\vspace{0.15em}
  \noindent\textit{\textbf{Answer:}~#3}
  \par
}

% ===== Skipped Question Note Macro =====
% Usage: \SkipBlock{Q1}{reason}

\newcommand{\SkipBlock}[2]{%
  \noindent\textcolor{SkipGray}{\textit{[#1 omitted: #2]}}%
  \vspace{0.5em}\par
}

% ===== Document Start =====

\begin{document}

\begin{center}
{\huge \textbf{Physics: Quantum Mechanics}}\\[0.25em]
{\Large Isomorphic Practice Questions (Wave Functions and Infinite Square Well)}\\[0.25em]
{\large Questions generated below. Multiple-choice answers in RED; incorrect answers in BLUE.}
\end{center}

\vspace{0.75em}

\MPBlock{Q1}{A particle of mass \(m\) in one dimension has the wave function
\[
\Psi(x,t)=C\,\exp\!\left[-\alpha\left(\frac{m x^2}{\hbar}\right)\right]\,e^{-i\Omega t},
\]
where \(\alpha>0\) and \(\Omega\) are real constants.}{
  \OEPart{a}{Find the normalization constant \(C\) (take \(C\) real and positive).}{Since the global phase does not affect probabilities, \(|\Psi(x,t)|^2=C^2\exp\!\left[-2\alpha\left(\frac{m x^2}{\hbar}\right)\right]\). Normalization requires \(\int_{-\infty}^{\infty} |\Psi|^2\,dx=1\), so
\[
C^2\int_{-\infty}^{\infty}e^{-\frac{2\alpha m}{\hbar}x^2}\,dx= C^2\sqrt{\frac{\pi\hbar}{2\alpha m}}=1,
\]
which gives \(C=\left(\frac{2\alpha m}{\pi\hbar}\right)^{1/4}.\)}
  \OEPart{b}{Write the integral expression used to compute \(\langle x\rangle\) and evaluate it for this state.}{Use the Born-rule expectation value
\[\langle x\rangle=\int_{-\infty}^{\infty}\Psi^*(x,t)\,x\,\Psi(x,t)\,dx.\]
Here \(|\Psi|^2=C^2 e^{-\frac{2\alpha m}{\hbar}x^2}\) is an even function of \(x\), so the integrand \(x|\Psi|^2\) is odd and therefore \(\langle x\rangle=0\).}
}

\OEBlock{Q2}{At time \(t=0\), a particle is described by the (real) wave function
\[
\Psi(x,0)=\begin{cases}
A\,(x/c), & 0\le x\le c\\
A\,(L-x)/(L-c), & c\le x\le L\\
0, & \text{otherwise}
\end{cases}
\]
where \(0<c<L\) and \(A=\dfrac{5}{L}\). What is the probability of finding the particle in the region \(x<c\)?}{The probability to the left of \(c\) is \(P(x<c)=\int_{-\infty}^{c}|\Psi(x,0)|^2dx=\int_{0}^{c}\left(A\frac{x}{c}\right)^2dx\). Evaluating gives \(P=\frac{A^2}{c^2}\cdot\frac{c^3}{3}=\frac{A^2 c}{3}\). With \(A=\frac{5}{L}\),
\[P(x<c)=\frac{25c}{3L^2}.\]}

\OEBlock{Q3}{A state is given by
\[
\Phi(x,t)=\varphi(x)\,e^{-i\varepsilon t/\hbar},
\]
where \(\varepsilon\) is a constant. Using the probability density, explain why this is a stationary state.}{The probability density is
\[|\Phi(x,t)|^2=\Phi^*(x,t)\Phi(x,t)=|\varphi(x)|^2\,\big(e^{+i\varepsilon t/\hbar}e^{-i\varepsilon t/\hbar}\big)=|\varphi(x)|^2,\]
which is independent of time because the phase factors cancel. Since all position probabilities are time-independent, the state is stationary.}

\MPBlock{Q4}{A normalized wave function can be expanded in a complete set of Hamiltonian energy eigenstates \(\{u_n(x)\}\) as
\[
\Psi(x,t)=\sum_{n=1}^{\infty} d_n\,u_n(x)\,e^{-i\mathcal{E}_n t/\hbar}.
\]
Assume the \(u_n\) are orthonormal and satisfy \(\hat H u_n=\mathcal{E}_n u_n\).}{
  \OEPart{a}{What is the physical meaning of \(|d_n|^2\)?}{In an energy eigenbasis, \(|d_n|^2\) is the probability that an energy measurement yields the eigenvalue \(\mathcal{E}_n\) (equivalently, that the post-measurement state is found in eigenstate \(u_n\)). For a normalized state, the probabilities satisfy \(\sum_{n=1}^{\infty}|d_n|^2=1\).}
  \OEPart{b}{Write the expectation value \(\langle H\rangle\) in terms of \(|d_n|^2\) and \(\mathcal{E}_n\).}{Because the measurement outcomes are \(\mathcal{E}_n\) with probabilities \(|d_n|^2\), the energy expectation value is the probability-weighted average
\[\langle H\rangle=\sum_{n=1}^{\infty}|d_n|^2\,\mathcal{E}_n.\]}
}

\MPBlock{Q5}{Consider a one-dimensional infinite square well of width \(L\):
\[
V(x)=\begin{cases}
0, & 0<x<L\\
\infty, & \text{otherwise.}
\end{cases}
\]
Inside the well define \(k^2\equiv 2mE/\hbar^2\).}{
  \OEPart{a}{Write the time-independent Schr\"odinger equation in the region \(0<x<L\) and give its general solution there.}{For \(0<x<L\), \(V=0\) so
\[-\frac{\hbar^2}{2m}\frac{d^2\psi}{dx^2}=E\psi \quad\Rightarrow\quad \psi''+k^2\psi=0\ \ (k^2=2mE/\hbar^2).\]
The general solution inside the well is \(\psi(x)=A\sin(kx)+B\cos(kx)\), and \(\psi=0\) outside because the potential is infinite.}
  \OEPart{b}{Apply the boundary conditions to find the allowed energies \(E_n\).}{The infinite walls impose \(\psi(0)=0\Rightarrow B=0\) and \(\psi(L)=0\Rightarrow \sin(kL)=0\). Thus \(kL=n\pi\) with \(n=1,2,3,\dots\), so \(k_n=\frac{n\pi}{L}\) and
\[E_n=\frac{\hbar^2 k_n^2}{2m}=\frac{n^2\pi^2\hbar^2}{2mL^2}.\]}
  \OEPart{c}{Find the normalized eigenfunctions \(\psi_n(x)\) in the interval \(0<x<L\).}{With \(B=0\) and \(k_n=n\pi/L\), the eigenfunctions are \(\psi_n(x)=A_n\sin\left(\frac{n\pi x}{L}\right)\) for \(0<x<L\). Normalization requires \(\int_0^L |\psi_n(x)|^2dx=1\), giving \(A_n=\sqrt{\frac{2}{L}}\). Hence
\[\psi_n(x)=\sqrt{\frac{2}{L}}\sin\left(\frac{n\pi x}{L}\right)\quad (0<x<L),\]
and \(\psi_n=0\) outside.}
  \OEPart{d}{Show that \(\psi_n\) and \(\psi_m\) are orthogonal when \(n\ne m\).}{Using the normalized functions,
\[\int_0^L \psi_n(x)\psi_m(x)\,dx=\frac{2}{L}\int_0^L \sin\left(\frac{n\pi x}{L}\right)\sin\left(\frac{m\pi x}{L}\right)dx.\]
The sine functions form an orthogonal set on \([0,L]\), so this integral equals \(0\) for \(n\ne m\) (e.g., by using \(\sin\alpha\sin\beta=\tfrac12[\cos(\alpha-\beta)-\cos(\alpha+\beta)]\) and evaluating at the limits).}
}

\OEBlock{Q6}{(Optional) A particle in a 1D infinite square well of width \(L\) is prepared at \(t=0\) in an equal superposition of the \(n=2\) and \(n=5\) stationary states:
\[
\Psi(x,0)=\frac{1}{\sqrt{2}}\big(\psi_2(x)+\psi_5(x)\big).
\]
Find the expectation value of the energy \(\langle E\rangle\).}{For \(\Psi=\frac{1}{\sqrt{2}}(\psi_2+\psi_5)\), the expansion coefficients have magnitudes \(|c_2|^2=|c_5|^2=\tfrac12\), so
\[\langle E\rangle=\tfrac12 E_2+\tfrac12 E_5=\frac{E_2+E_5}{2}.\]
In an infinite well, \(E_n=\frac{n^2\pi^2\hbar^2}{2mL^2}\), hence
\[\langle E\rangle=\frac{1}{2}\left(\frac{4\pi^2\hbar^2}{2mL^2}+\frac{25\pi^2\hbar^2}{2mL^2}\right)=\frac{29\pi^2\hbar^2}{4mL^2}.\]}

\end{document}
